We present BEAMFS v2, a Linux filesystem providing autonomic
recovery from electromagnetic perturbations of persistent storage,
including stochastic single-event upsets, multi-bit upsets, and
adversarial high-power electromagnetic events. BEAMFS v2 follows
the FTRFS lineage of Fuchs et al.~\cite{fuchs2015ftrfs} and
extends it with two contributions:
\textbf{(i)}~an INLINE protection scheme applying Reed--Solomon
\(\mathrm{RS}(255,239)\) per 255-byte subblock with sixteen
subblocks per 4\,KiB disk block, recovering up to eight symbol
errors per subblock and up to one hundred and twenty-eight symbol
errors per block, on the read path, in place, and persisted to
disk byte-perfect (\emph{autonomic in-place repair});
\textbf{(ii)}~an explicit \emph{saturation observability} contract:
when a perturbation exceeds the per-subblock correction capacity,
the recovery operator returns fail-closed and emits a journal entry
flagged uncorrectable, making the saturation event observable
post-mortem rather than silent.

The threat model is restated in electromagnetic terms (\emph{EMR
threat model}) covering both stochastic perturbations
(Family\,A, modelled by canonical SEU/MBU rate calculations) and
adversarial perturbations (Family\,B, IEMI/HPM/EMP regimes per
IEC\,61000-2-13 and 61000-4-36), bounded by a transverse saturation
boundary at the per-subblock correction capacity.

The recovery soundness theorem of the v1 report (Theorem\,IV.1)
was empirically falsified by the companion fault injection
operator RadFI~\cite{desbrieres2026radfi} on 2026-04-28. The
present report retracts that theorem, reformulates the recovery
calculus, and presents two replacement statements:
\textbf{Theorem v2.1} (recovery under bounded EMR perturbation),
which retains the termination/safety/coverage structure of
Theorem\,IV.1 with an EMR-stated hypothesis on the perturbation
family; and \textbf{Theorem v2.2} (saturation observability),
which formalizes the contract that fail-closed recovery emits an
auditable uncorrectable-flagged journal entry at the affected
\((\text{block},\text{subblock})\) coordinates.

Empirical results are reported on a Yocto-built aarch64 research
image: byte-level disk corruption of an INLINE-formatted volume is
recovered byte-perfect (\emph{Stage 3}, recovery validation
milestone); read latency on the
INODE\_UNIVERSAL protection scheme matches \texttt{ext4} on median
and 95th percentile, with a 1.85\(\times\) p99 tail (\emph{Stage 4 / B1}, performance
evaluation milestone); write+fsync latency carries an
\(\sim\)8\(\times\) median
overhead versus \texttt{ext4} due to per-block RS encoding, with a
tighter overall distribution (1.74\(\times\) p50/p99 spread for
BEAMFS versus 5.00\(\times\) for \texttt{ext4}, Stage 4 / B2);
INLINE read on the conformance fixture canary block, ten runs of
\(\sim\)38{,}000 reads each, exhibits zero RS correction events
and a stable fixture SHA-256 (Stage 4 / B3). The kernel module
totals 2{,}198 lines of C, well under the 5{,}000-line audit
budget targeted for safety-critical certification regimes.

The reference implementation is publicly available at
\texttt{roastercode/beamfs}~\cite{beamfs-impl}, GPG-signed, with
all empirical results reproducible from the tagged commit.
